% ---------------------------------------------------------------------------
% Hardware sections for the marble-maze CPS report.
% Style/structure mirrors the ConnecTUM (Connect 4) CPS report:
%   §Mechanical design  →  requirements → design → motivation → integration
%   §Electrical and control system → architecture, per-actuator control,
%   wiring table, iteration/failure narratives.
% Figures expected in fig/. [CONFIRM: ...] marks numbers to verify.
% ---------------------------------------------------------------------------

\section{System overview}
\label{sec:overview}

The system autonomously solves a physical wooden labyrinth by tilting the
board about two orthogonal axes so that a steel marble follows a target path
from start to goal. It is a cyber-physical system composed of four subsystems:
(i) a commercial wooden labyrinth whose inner gimbal is tilted about two axes;
(ii) a two-servo actuation stage that drives those axes; (iii) an overhead
global-shutter camera that provides state feedback; and (iv) a host PC that
runs perception and control in Python and issues tilt commands over a USB
serial link to an Arduino microcontroller. The Arduino translates high-level
tilt commands into pulse-width signals through a dedicated PWM driver.

Following the same philosophy as the mechanical iterations described below, we
built the platform incrementally: we first brought up the electronics and
firmware on the bench with the servos unloaded, validated the command
interface and safety behaviour, and only then progressed to mechanical
coupling, camera mounting, and calibration. This staged approach let us
isolate electrical, mechanical, and software faults independently, which
proved valuable when diagnosing the failures reported in
Sections~\ref{sec:mech} and \ref{sec:elec}.

\section{Mechanical design}
\label{sec:mech}

This section describes the mechanical components, their design, and their
integration. Unlike a purpose-built machine, our starting point is an
unmodified commercial labyrinth; the mechanical work therefore centres on
\emph{interfacing} servo actuators to the board's existing tilt mechanism
without compromising its motion. We favoured reversible, low-part-count
solutions and iterated over several coupling concepts, keeping the design that
best balanced simplicity, backlash, and safe travel.

\subsection{Board and tilt mechanism}
\label{sec:mech-board}

The labyrinth is a two-axis gimbal: an outer frame pivots about one axis and an
inner frame about the orthogonal axis, giving two independent tilt degrees of
freedom. Natively, each axis is driven by a knob coupled through a
string-and-spring transmission. The play area is approximately
$335\times335$\,mm [CONFIRM: measured board dimensions]. Two properties of this
mechanism shaped our design: the string/spring transmission introduces
compliance and a limited, spring-loaded travel range, and the two axes have
hard mechanical end-stops that must not be driven into under power.

\subsection{Actuation requirements}
\label{sec:mech-req}

For the actuation interface we formulated the following requirements:

\begin{enumerate}
  \item Drive each of the two tilt axes independently across the range needed
    for marble control.
  \item Keep commanded travel strictly inside the mechanism's mechanical
    end-stops to avoid stalling the servos.
  \item Minimise backlash between the servo and the board, since slop degrades
    the accuracy of open-loop tilt commands.
  \item Avoid irreversible modification of the board where practical.
  \item Provide a repeatable neutral (level) position.
\end{enumerate}

We selected two MG996R digital metal-gear servos (13\,kg$\cdot$cm class)
[CONFIRM: exact model/torque], one per axis. A key early observation was that
the marble is controlled by only a few degrees of \emph{board} tilt, which
corresponds to a modest actuator swing; full travel of the native mechanism is
not required. This relaxed the range requirement and made hobby servos viable.

\subsection{Actuator coupling: from pushrod linkage to direct drive}
\label{sec:mech-linkage}

The coupling between servo and board went through two iterations.

\paragraph{Iteration~A --- ball-joint pushrod (superseded).}
Our first concept followed a conventional RC pushrod pattern
(Figure~\ref{fig:tie-rod}): a ball-joint rod end threaded onto an M3$\times$100
threaded rod at each end, connecting an offset servo horn to a lever fixed on
the board's knob. A hole drilled off-centre in the knob set the lever length
and hence the servo-to-axis gear ratio. This approach is forgiving of servo
placement, but in our setup it introduced backlash at the two ball joints and a
nonlinear angle map that complicated calibration, and it required a precisely
located lever hole to match the servo's limited swing. These drawbacks
motivated a change of approach.

\paragraph{Iteration~B --- direct coupling to the tilt axes (adopted).}
In the final design, each servo is mounted beneath the board and its output is
coupled directly to a gimbal tilt axis, driving the inner and outer frames
respectively (Figure~\ref{fig:direct-coupling}) [CONFIRM: exact coupling ---
set-screw hub / printed adapter / horn bolted to pivot]. Direct coupling
removes the ball-joint backlash and yields a near 1:1 servo-to-axis angle map,
simplifying the tilt command model at the cost of requiring accurate coaxial
alignment of the servo to the pivot. Because the mapping is direct, the servo
travel must be constrained so it cannot drive the axis into its mechanical
stop; we enforce this in firmware (Section~\ref{sec:elec-safety}).

\subsection{Servo mounting and camera support}
\label{sec:mech-mount}

Both servos are fixed rigidly relative to the board so that reaction torque
produces board tilt rather than actuator displacement; any compliance in the
servo mount would appear directly as tilt error. The overhead camera is
supported on an articulated arm and clamp and is referenced to the same rigid
base as the board, so that incidental disturbances move the camera and board
together and preserve the image-to-board calibration
(Section~\ref{sec:calibration}).

\begin{figure}[t]
\centering
% \includegraphics[width=\linewidth]{fig/tie_rod_linkage.jpg}
\caption{Iteration~A: ball-joint pushrod linkage concept, in which a threaded
rod with ball-joint ends connects an offset servo horn to a lever on the board
knob. Illustrative reference; superseded in our build by direct coupling.}
\label{fig:tie-rod}
\end{figure}

\begin{figure}[t]
\centering
% \includegraphics[width=\linewidth]{fig/direct_coupling.jpg}
\caption{Iteration~B (adopted): the servo output is coupled directly to the
board's tilt pivot, eliminating rod-end backlash.}
\label{fig:direct-coupling}
\end{figure}

\section{Electrical and control system}
\label{sec:elec}

This section describes the electronics that drive the servos and the firmware
that mediates between the host PC and the actuators. As with the mechanical
design, the final architecture is the result of iterations, two of which
(power distribution and the serial/firmware interface) are described in detail
because they materially affected reliability.

\subsection{Architecture}
\label{sec:elec-arch}

We formulated the following requirements for the control electronics:

\begin{enumerate}
  \item Sufficient host compute for real-time perception on the camera stream.
  \item Deterministic, low-latency generation of servo PWM signals.
  \item A power scheme able to supply servo stall currents without disturbing
    the logic and camera electronics.
  \item A simple, robust host--microcontroller command interface.
\end{enumerate}

Perception and control run on a host PC (equipped with an NVIDIA RTX~4050)
[CONFIRM: CPU/RAM/OS], which communicates over USB serial with an Arduino
UNO~R4 Minima. The Arduino does not generate PWM directly; instead it drives a
PCA9685 16-channel PWM controller over I\textsuperscript{2}C
(address~0x40, 50\,Hz servo frame), with channel~0 assigned to the yaw axis and
channel~1 to the pitch axis. Offloading PWM generation to the PCA9685 keeps
timing off the microcontroller's main loop and matches the pattern of using a
dedicated PWM driver behind a microcontroller. This split (host for
high-level logic, microcontroller + PWM driver for real-time actuation)
mirrors the distributed, serial-connected architecture we found effective for
decoupling flexible Python logic from real-time motor control.

\subsection{Power distribution and stall protection}
\label{sec:elec-power}

Servo power is supplied by a dedicated 5\,V / 10\,A DC supply wired to the
PCA9685 V+ terminal, electrically separate from the Arduino logic supply, which
is delivered over USB. The two domains share a common ground at the PWM driver.
This split-rail arrangement isolates inductive servo transients from the logic
and camera electronics.

An early failure made the importance of this subsystem concrete. During an
actuation test the servo was commanded past the mechanism's mechanical stop;
unable to reach the setpoint, it drew sustained stall current
([CONFIRM: $\sim$2.5\,A per MG996R]). The initial harness used thin jumper wire
(\textasciitilde26\,AWG) on the servo rail, which cannot carry this current and
overheated to the point of conductor failure. We resolved this in two ways:
(i) the servo-rail conductors were replaced with 18--20\,AWG wire and an inline
fuse was added, so that an over-current condition blows a fuse rather than a
conductor; and (ii) commanded travel was constrained in firmware
(Section~\ref{sec:elec-safety}) so the actuator cannot be driven into a stall
in the first place. We report this because undersized power wiring and
unconstrained travel are common, silent hazards in servo-actuated rigs.

\begin{table}[t]
\centering
\small
\begin{tabular}{@{}ll@{}}
\toprule
\textbf{Signal} & \textbf{Connection} \\
\midrule
\multicolumn{2}{@{}l}{\textit{Arduino UNO R4 Minima}} \\
\quad 5V   & PCA9685 VCC (logic) \\
\quad GND  & PCA9685 GND (common ground) \\
\quad SDA  & PCA9685 SDA \\
\quad SCL  & PCA9685 SCL \\
\quad USB  & Host PC (serial) \\
\midrule
\multicolumn{2}{@{}l}{\textit{PCA9685 PWM driver}} \\
\quad V+   & 5\,V / 10\,A supply (+), via fuse \\
\quad GND  & 5\,V / 10\,A supply ($-$) \\
\quad CH0  & Yaw servo (signal/V+/GND) \\
\quad CH1  & Pitch servo (signal/V+/GND) \\
\bottomrule
\end{tabular}
\caption{Wiring of the control electronics. Logic power (Arduino, USB) and
servo power (5\,V/10\,A supply) are separate rails sharing a common ground at
the PCA9685.}
\label{tab:wiring}
\end{table}

\subsection{Firmware and command interface}
\label{sec:elec-firmware}

The Arduino runs a minimal firmware exposing a line-oriented serial protocol.
The host sends normalized tilt commands \texttt{SET <yaw> <pitch>}, each value
in $[-1, 1]$, which the firmware maps to a servo pulse width about a
1500\,\textmu s neutral. The protocol also provides \texttt{PING}/\texttt{PONG}
liveness checks and an explicit \texttt{NEUTRAL} command; malformed input is
rejected with an error response.

The serial/firmware interface itself was tuned over two iterations, motivated
by control-loop latency:

\begin{itemize}
  \item \textbf{Throughput.} The serial link was raised from 115\,200 to
    500\,000\,baud, and the I\textsuperscript{2}C clock to the PCA9685 to
    400\,kHz (Fast Mode), reducing the per-command transport time
    ([CONFIRM: $\sim$1.7\,ms $\rightarrow$ $\sim$0.3\,ms per command]).
  \item \textbf{Loop structure.} The initial firmware read a full line with a
    blocking call and echoed an acknowledgement per command. We replaced this
    with a non-blocking, character-by-character reader and dropped the
    per-command echo, so the main loop never stalls on serial I/O. Servo pulses
    are refreshed on a fixed 200\,Hz schedule independent of command arrival.
\end{itemize}

\subsection{On-board safety behaviour}
\label{sec:elec-safety}

Three safety mechanisms run on the microcontroller, independent of the host, so
that a software fault, a disconnected cable, or an out-of-range command cannot
damage the mechanism or leave the board tilted:

\begin{enumerate}
  \item \textbf{Travel clamp.} Commanded pulse widths are hard-clamped to a
    safe band about neutral ($1500\pm300$\,\textmu s in the current
    configuration), keeping the actuator inside the mechanism's end-stops.
  \item \textbf{Slew limiting.} Each update ramps the commanded pulse toward
    its target rather than stepping instantaneously, avoiding abrupt torque
    transients.
  \item \textbf{Watchdog.} If no host command is received within 500\,ms, both
    axes are returned to neutral, so a lost connection fails safe to a level
    board.
\end{enumerate}

The travel clamp is deliberately conservative: after the stall incident of
Section~\ref{sec:elec-power}, we evaluated widening it to increase usable tilt
range but reverted to the conservative band, prioritising actuator safety over
maximum travel.

\subsection{Sensing}
\label{sec:elec-sensing}

State feedback is provided by an overhead USB2.0 UVC global-shutter monochrome
camera (OV9281-class sensor, 720p/120\,fps advertised) [CONFIRM: configured
resolution/FPS]. A global shutter avoids the motion-induced skew that a rolling
shutter would introduce on the moving marble. Camera intrinsics and the
image-to-board homography are estimated in a calibration step
(Section~\ref{sec:calibration}).

% ---------------------------------------------------------------------------
% Open placeholders for the remaining (software) sections.
% ---------------------------------------------------------------------------

\section{Calibration}
\label{sec:calibration}
% TODO (software): camera intrinsics; image-to-board homography; procedure and
% residual error.

\section{Perception}
\label{sec:perception}
% TODO (software): marble detection/tracking; state (position, velocity)
% estimation; lost-ball and hole-fall handling.

\section{Planning and control}
\label{sec:control}
% TODO (software): path annotation; path-following controller; command
% saturation and jerk limiting; per-section tuning.

\section{Evaluation}
\label{sec:evaluation}
% TODO: success rate, completion time, and failure modes across repeated runs.
